\clearpage\begingroup\NoSourceHead\SourceLeaf{208}
\newcommand{\ChronMoon}{\mathord{\vcenter{\hbox{\begin{tikzpicture}[x=1ex,y=1ex]
\path[use as bounding box] (0,0) rectangle (.65,1.1);
\path[fill=black] (.6,1.1) .. controls (-.1,1.05) and (-.1,.05) .. (.6,0)
.. controls (.13,.2) and (.13,.9) .. (.6,1.1) -- cycle;
\end{tikzpicture}}}}}
\begin{center}
{[III.]}\\[.8em]
\textbf{[DE CALENDARIO ECCLESIASTICO.]}\\[.7em]
\rule{\textwidth}{.35pt}\\[.2em]
{\small P. \textsc{Tittel}, Methodus technica, brevis, perfacilis ac perpetua construendi Calendarium ecclesiasticum stylo\\
tam novo quam vetere, Goettingae 1816, pag. 7, 11, 13.}\\[.2em]
\rule{\textwidth}{.35pt}\\[1em]
\{§ 8.\\[.7em]
\textbf{Problema.}\\[.7em]
\textls[100]{Determinare numerum festivalem Stylo veteri pro quocunque anno dato.}\\[.7em]
\textbf{Solutio III.}
\end{center}
\noindent\hspace*{\parindent}Indice anni dati successive per numeros 19, 4, 7 diviso, oriatur ex divisione
\begin{center}
\begin{tabular}{@{}r l l r l l@{}}
1)&indicis dati . . .&per&19,&residuum&$a$\\
2)&indicis dati . . .&per&4,&residuum&$b$\\
3)&indicis dati . . .&per&7,&residuum&$c$\\
4)&numeri $19a+15$ . . .&per&30,&residuum&$d$\\
5)&numeri $2b+4c+6d+6$ . . .&per&7,&residuum&$e$.
\end{tabular}
\end{center}
\noindent\hspace*{\parindent}Erit absque ulla exceptione numerus festivalis quaesitus, aequalis summae: $d+e+1$.
\begin{center}
§ 12.\\[.7em]
\textbf{Problema.}\\[.7em]
\textls[100]{Calculare aequationem lunaesolarem pro quocunque anno assignato.}\\[.7em]
\textbf{Solutio.}
\end{center}
\noindent\hspace*{\parindent}Subtrahatur aequatio lunaris anni assignati per (§ 11.) calculata ab aequatione solari eiusdem anni, per (§ 10.) definita; differentia dividatur per 30,
\HistoricalFoot{}{26\textsuperscript{*}}
\clearpage\SourceLeaf{209}\HistoricalHead{204}{L}{CHRONOLOGIE.\quad VERÖFFENTLICHUNGEN.}
\noindent neglectoque quotiente notetur solummodo residuum, istud ipsum erit aequatio lunaesolaris quaesita. Si itaque aequationem solarem anni assignati $\odot$, lunarem $\ChronMoon$, lunaesolarem $\odot\ChronMoon$ vocemus; erit $\odot\ChronMoon={}$ residuo numeri $(\odot-\ChronMoon)$, per 30 divisi.
\begin{center}\textbf{Exempla.}\end{center}
\noindent\hspace*{\parindent}1) Aequatio lunaesolaris anni 1815 est: 6. Erat enim pro hoc anno in (§ 10.) $\odot=14$, in (§ 11.) $\ChronMoon=8$, adeoque $\odot\ChronMoon=6$.

2) Aequatio lunaesolaris anni 1783 erat: aeque 6. Erat enim pro anno hocce in (§ 10.) $\odot=13$, in (§ 11.) $\ChronMoon=7$, unde $\odot\ChronMoon=6$ sequitur.

3) Aequatio lunaesolaris anni 28066 erit 28. Nam pro hoc anno obtinuimus in (§ 10.) $\odot=210$, in (§ 11.) $\ChronMoon=92$, hinc $\odot-\ChronMoon=118$, adeoque $\odot\ChronMoon=28$.
\begin{center}
§ 14.\\[.7em]
\textbf{Problema.}\\[.7em]
\textls[100]{Determinare numerum festivalem cuiuscunque anni dati stylo novo.}\\[.7em]
\textbf{Solutio II.}
\end{center}
\noindent\hspace*{\parindent}A. Oriatur ex divisione:
\begin{center}
\begin{tabular}{@{}r l l r l l@{}}
1)&indicis anni dati&per&19,&residuum&$a$\\
2)&indicis anni dati&per&4,&residuum&$b$\\
3)&indicis anni dati&per&7,&residuum&$c$\\
4)&numeri $19a+17+\odot\ChronMoon$&per&30,&residuum&$d$\\
5)&numeri $2b+4c+6d+4+\odot$&per&7,&residuum&$e$.
\end{tabular}
\end{center}
\noindent\hspace*{\parindent}B. Numerus festivalis quaesitus erit aequalis numero: $(d+e+1)$, exceptis duobus casibus:

1) Si $d+e+1=36$ inveniatur. Aut

2) Si fuerit $d+e+1=35$, simulque $d=28$, $e=6$, et aequatio lunaesolaris aliquem ex sequentibus valoribus obtinuerit: 4, 7, 12, 15, 18, 23, 26, 29.

His enim duobus casibus numerus festivalis, lege ordinaria obtentus 7 unitatibus mulctandus est, adeoque fiet numerus festivalis casu primo $36-7=29$; casu vero secundo $28=35-7$.\}
\clearpage\SourceLeaf{210}\HistoricalHead{205}{R}{BEMERKUNGEN.}
\begin{center}\textbf{BEMERKUNGEN.}\end{center}
\begingroup\small
\noindent\hspace*{\parindent}Die vorstehend abgedruckten Stellen aus \textsc{Tittel}s Werk werden durch die Worte der Einleitung: »Solutiones elegantissimas: in 8. § tertiam, in 12. § primam, in 14. § secundam in acceptis referendas habeo Viro celeberrimo \textsc{Carolo Friderico Gauss}« als von \textsc{Gauss} herrührend gekennzeichnet.

Da März $21+{}$ Festzahl das Osterdatum ergibt, so stimmt für den julianischen Kalender die Solutio III des § 8. mit der \textsc{Gauss}schen Vorschrift in den Werken VI, S. 78 völlig überein.

Nach \textsc{Tittel} (\textit{Methodus technica} § 11., S. 10) lautet die Mondgleichung
\[
\ChronMoon=Q\left(\frac{13+8k}{25}\right)+2
\]
\noindent und (ebenda § 10., S. 9) die Sonnengleichung
\[
\odot=k-Q\left(\frac{k}{4}\right)^{*)},
\]
\begingroup\renewcommand{\thefootnote}{}\footnotetext{*) $k$ bedeute die Säkularzahl des Jahres $A$, also $k=Q\left(\frac{A}{100}\right)$; $Q$ und $R$ haben die Bedeutung von Quotient und Rest.}\endgroup
\noindent also ist (siehe oben § 12.)
\[
\odot\ChronMoon=R\left(\frac{k-Q\left(\dfrac{k}{4}\right)-Q\left(\dfrac{13+8k}{25}\right)-2}{30}\right);
\]
\noindent mithin ergibt sich für die von \textsc{Gauss} (Werke VI, S. 78) mit $M$ bezeichnete Größe, wenn man gemäß der \textit{Berichtigung} (oben S. 201) $p=Q\left(\frac{13+8k}{25}\right)$ nimmt, die Kongruenz $M\equiv\odot\ChronMoon+17$ (mod. 30), und weiter wird $N\equiv\odot+4$ (mod. 7). Infolgedessen tritt für \textsc{Tittel}s Größe 4) der Ausdruck $19a+M$ und für seine Größe 5) der Ausdruck $2b+4c+6d+N$. Dies stimmt mit der von \textsc{Gauss} für den gregorianischen Kalender gegebenen Vorschrift überein (Werke VI, S. 78), wenn man bedenkt, daß Ostern an dem durch März $21+{}$ Festzahl bestimmten Tage gefeiert wird.
\begin{flushright}\textsc{Loewy.}\end{flushright}
\endgroup
\begin{center}\rule{.24\textwidth}{.35pt}\end{center}
\clearpage\endgroup
